\section{Conclusion}\label{sec:conclusion}

Theorem~\ref{thm:main} answers Question~46 of \cite{BBBHHL}, repeated
as Problem~10 of \cite{BrechBrendleTelles}, by showing that $\mathfrak{rr}=\nonM$.
In this section, we record some further
consequences for other related cardinal characteristics of the continuum.

In \cite{BBBHHL}, the authors defined several variants of the rearrangement number
by restricting the permitted outcomes of rearrangements.

\begin{definition}\label{def:rr-variants}
If $\pi\in\Sym$
rearranges $a\in\CC$, we say that the outcome of the rearrangement is of type
\begin{enumerate}
\item[$f$] if $\sum_n a_{\pi(n)}$ converges to a real number different
      from the original sum;
\item[$i$] if it diverges to $+\infty$ or to $-\infty$;
\item[$o$] if its partial sums diverge by oscillation, meaning that
      they have no limit in the extended real line
      $[-\infty,+\infty]$.
\end{enumerate}
If $x$ is one of $f,i,o,fi,fo,io$, a collection
$\Pi\subseteq\Sym$ is \emph{$x$-rearranging} if, for every
$a\in\CC$, some $\pi\in\Pi$ has one of the outcome types whose letter
occurs in $x$. Define $\mathfrak{rr}_x$ as the least cardinality of an
$x$-rearranging family. This gives the cardinals
\[
 \rrf,\quad\rri,\quad\rro,\quad
 \rrfi,\quad\rrfo,\quad\rrio.
\]
\end{definition}

Some ZFC relations between these cardinals were obtained in \cite{BBBHHL} and are shown in Figure~\ref{fig:rr-variants}.

\begin{figure}[ht]
\begin{center}
\begin{tikzpicture}[
  x=1cm,
  y=1cm,
  cardinal/.style={font=\large,inner sep=2pt}
]
  \node[cardinal] (rrf)  at (-1.5,1.5) {$\rrf$};
  \node[cardinal] (rri)  at ( 1.5,1.5) {$\rri$};
  \node[cardinal] (rrfi) at ( 0,0)     {$\rrfi$};
  \node[cardinal] (rr)   at ( 0,-1.5)  {$\rr=\rrfo=\rrio=\rro$};

  \draw (rrf) -- (rrfi);
  \draw (rri) -- (rrfi);
  \draw (rrfi) -- (rr);
\end{tikzpicture}
\end{center}
\caption{The hierarchy of rearrangement numbers.
Higher means provably at least as large in ZFC.}
\label{fig:rr-variants}
\end{figure}

Recall that $\dd$ is the least size of a dominating family in
$(\omega^\omega,\leq^*)$, and $\cofM$ is the least size of a family
of meager sets cofinal under inclusion in the ideal $\cM$.

\begin{corollary}\label{cor:rr-variants}
In ZFC,
\[
 \rr=\rro=\rrfo=\rrio=\nonM
 \quad\text{and}\quad
 \cofM\leq\rrfi\leq\min\{\rrf,\rri\}.
\]
In particular, if $\cofM=\mathfrak c$, then
$\rrfi=\rrf=\rri=\mathfrak c$.
\end{corollary}

\begin{proof}
By \cite[Theorem~5]{BBBHHL}, $\rr=\rro=\rrfo=\rrio$, so the first
equality follows from Theorem~\ref{thm:main}. Also,
$\dd\leq\rrfi$ by \cite[Theorem~12]{BBBHHL}. Thus the classical
identity $\cofM=\max\{\dd,\nonM\}$
\cite[Theorem~5.6]{BlassHandbook} gives
\[
 \cofM=\max\{\dd,\rr\}\leq\rrfi.
\]
The remaining conclusions follow from the definitions and the upper
bound $\mathfrak c$ for all the rearrangement numbers.
\end{proof}

Some cardinals related to the rearrangement number were defined
in \cite{BrendleBrianHamkins}, including the subseries number $\sz$,
the subrearrangement number $\sr$, and the oscillation subseries number $\sz_o$.

\begin{definition}[{\cite{BrendleBrianHamkins}}]\label{def:subseries}
The \emph{subseries number} $\sz$ is the least size of a family
$\mathcal A\subseteq[\omega]^\omega$ such that, for every $a\in\CC$ there exists $A\in\mathcal A$ such that $\sum_{n\in A}a_n$ diverges.
Restricting divergence to divergence by oscillation gives the \emph{oscillation subseries number} $\sz_o$.

The \emph{subrearrangement number} $\sr$ is the least size of a
family $F$ of injections from $\omega$ to $\omega$ such that, for
every $a\in\CC$, some $f\in F$ makes $\sum_n a_{f(n)}$ diverge.
\end{definition}

Brendle, Brian, and Hamkins proved
$\covN\leq\sz\leq\sz_o\leq\nonM$
\cite[Theorems~4, 11, and~13]{BrendleBrianHamkins}.
They also established $\rr\leq\max\{\bb,\sz\}$ and
$\sz_o\leq\max\{\bb,\sz\}$
\cite[Theorems~19(1) and~21]{BrendleBrianHamkins}, while
$\bb\leq\rr$ was already known from \cite[Theorem~11]{BBBHHL}.
For subrearrangements, they obtained the identity
$\sr=\min\{\sz,\rr\}$
\cite[Theorem~35]{BrendleBrianHamkins}.

The comparison between $\sz$ and $\rr$ was left undecided in one
direction: Question~20 of \cite{BrendleBrianHamkins} asks whether
$\rr<\sz$ is consistent. Equivalently, in view of $\bb\leq\rr$,
they asked whether the bound $\rr\leq\max\{\bb,\sz\}$ can be strict.
At the end of Section~10, they also asked whether $\sr<\sz$ is
consistent. In the other direction, their analysis of the Laver model
already gave the consistency of $\sz=\sz_o<\bb=\rr$
\cite[Section~9]{BrendleBrianHamkins}.
Combining the known relations with Theorem~\ref{thm:main} yields the
following answers.

\begin{corollary}\label{cor:subseries}
In ZFC,
\[
 \sr=\sz\leq\sz_o\leq\rr=\nonM,
 \qquad
 \rr=\max\{\bb,\sz\}=\max\{\bb,\sz_o\}.
\]
If $\bb<\nonM$, then $\sr=\sz=\sz_o=\rr=\nonM$.
\end{corollary}

\begin{proof}
By \cite[Theorems~13 and~19(1)]{BrendleBrianHamkins},
$\sz_o\leq\nonM$ and $\rr\leq\max\{\bb,\sz\}$.
Together with $\sz\leq\sz_o$, $\bb\leq\rr$, and
Theorem~\ref{thm:main}, these yield
\[
 \rr\leq\max\{\bb,\sz\}\leq\max\{\bb,\sz_o\}\leq\rr.
\]
Moreover, \cite[Theorem~35]{BrendleBrianHamkins} gives
$\sr=\min\{\sz,\rr\}=\sz$.
If $\bb<\nonM=\rr$, the identity $\rr=\max\{\bb,\sz\}$ forces
$\sz=\rr$, and the other equalities follow.
\end{proof}

In particular, Question~20 of \cite{BrendleBrianHamkins}, asking
whether $\rr<\sz$ is consistent, has a negative answer. The same is
true of the question at the end of Section~10 of that paper about
the consistency of $\sr<\sz$.

\subsection*{Acknowledgments}

The author thanks Tan Özalp for pointing out a serious error concerning the use of
the Laver property in a forcing notion in an earlier version of this manuscript.

This study was financed, in part, by the São Paulo Research Foundation
(FAPESP), Brazil. Process Number 2025/07302-0.

\subsection*{Declaration on the use of generative artificial intelligence}

The author used the OpenAI GPT-6 Astra large language model through Codex during the
research and preparation of this manuscript for generating and
developing novel mathematical arguments and proofs, searching the literature,
checking arguments for possible gaps and errors, and drafting and
revising the text and \LaTeX{} source. The author
takes full responsibility for the arguments, claims, references, and
any errors that remain.
